% ============================================================
%  CHAPTER 1 --- INTRODUCTION
% ============================================================
\chapter{Introduction}
\label{chap:intro}

Labour Force Surveys (LFS) provide information on employment,
unemployment, and labour underutilisation. Their value depends on the
consistent application of statistical concepts and on the quality of the
responses collected. The International Labour Organization (ILO)
establishes the relevant conceptual framework through the International
Conference of Labour Statisticians~\cite{ilo2013workstats}. Within that
framework, classifying open descriptions of work is a measurement task:
the wording supplied by a respondent must be interpreted in relation to a
defined coding scheme.

This thesis investigates a UAE-oriented conversational LFS prototype.
The implemented questionnaire includes employment, education, demographic,
and working-condition items, together with UAE-specific fields such as
emirate and residence duration. The interface provides English, Arabic,
Urdu, Hindi, and Tagalog, with additional processing for Gulf Arabic.
The prototype integrates occupation coding to ISCO-08, industry coding to
ISIC Rev.~4, and education coding to ISCED~2011 and ISCED-F~2013. These
are the classification versions selected for this study; the prototype is
not an authorised national survey instrument.

The central problem is how to connect flexible language input with
reproducible classification and questionnaire control. Conversational
fluency alone does not establish correct coding, and successful software
execution does not establish valid measurement of a population. The
research therefore separates classification effectiveness, cross-standard
checks, language robustness, and operational feasibility.

\section{Motivation and Background}
\label{sec:motivation}

\subsection{Occupation Coding as a Measurement Problem}

ISCO-08 groups jobs according to the tasks and duties performed. Its
four-level structure contains 10 major groups, 43 sub-major groups,
130 minor groups, and 436 unit groups~\cite{ilo2012isco_v1}.
A description such as \emph{I repair equipment} may fit several unit groups
unless the equipment, tasks, or work setting is specified. A classifier
must therefore handle incomplete descriptions and distinguish related
occupations at different levels of detail.

Coding during an interview is an established research direction.
Schierholz et al.~\cite{schierholz2016interview} investigated supervised
candidate suggestions that respondents could select while an interview
was in progress. More recently, SOCbot combines retrieved occupation
candidates with LLM-based classification and dynamic follow-up
questions~\cite{kochar2025socbot}. These precedents motivate a study of
retrieval configurations and their integration into a broader survey;
they preclude treating real-time occupation coding or conversational
probing alone as a new contribution.

\subsection{Related but Distinct Classification Targets}

Occupation, industry, and education describe different aspects of a
respondent's situation. ISCO-08 concerns work tasks; ISIC classifies the
economic activity of the relevant establishment; ISCED describes education
programmes and attainment, while ISCED-F describes fields of education
and training~\cite{ilo2012isco_v1,unsd2008isic,unesco2012isced,unesco2015iscedf}.
For example, a software developer can work in a hospital without the
occupation and industry classifications being inconsistent. Similarly,
formal qualifications are not a complete description of the skills
acquired through experience.

Integrating these targets creates an opportunity for review of unusual
combinations, but it also creates a risk of rejecting legitimate cases.
This thesis consequently treats its Semantic Relation Engine (SRE) as a
heuristic screening mechanism. Its hand-built associations and severity
rules identify cases for review; they are not an official correspondence
table or proof that a respondent's answers are incorrect.

\subsection{Multilingual Input and Arabic Dialect Variation}

Arabic language processing involves Modern Standard Arabic, regional
varieties, and code-switching with other languages. Hamed et
al.~\cite{hamed2025codeswitching} review the resulting resource and
modelling challenges. In a survey, these challenges occur together with
specialised occupation terminology and short answers. Translating the
questionnaire is therefore only one part of providing language support:
answer interpretation, routing, and classification must also preserve the
meaning of each response.

The prototype uses multilingual embeddings and a dictionary of Gulf
Arabic normalisation rules. Its evaluation distinguishes the availability
of language-specific interface text from measured performance on labelled
examples. Neither multilingual model training nor a translated interface
establishes equivalent accuracy across the five intended languages.

\subsection{Questionnaire Control and Operational Feasibility}

Computer-assisted personal interviewing (CAPI) already supports automated
skip logic and validation. For example, World Bank Survey Solutions
evaluates enabling and validation conditions and supports multilingual
questionnaires~\cite{worldbank2016surveysolutionsvalidation,worldbank2026surveysolutionslanguages}.
The proposed system extends these familiar controls with interpretation
of conversational responses and classification during the same workflow.
Its contribution is evaluated as an integration of components rather
than as the invention of conditional routing or real-time validation.

The prototype also reuses stored responses through a Person Register,
records classification decisions, and provides human-in-the-loop (HITL)
review. These mechanisms can be tested for functional correctness.
Claims about reduced respondent burden, interview cost, or improved
participation require a comparison with human respondents and an
appropriate survey administration baseline. Such effects cannot be
inferred from question counts or automated processing time alone.

\begin{table}[htbp]
\centering
\caption{Research problems and the evidence required to assess them.}
\label{tab:ch1-challenges}
\small\renewcommand{\arraystretch}{1.3}
\begin{tabular}{>{\raggedright}p{3.0cm}
                >{\raggedright}p{5.0cm}
                >{\raggedright\arraybackslash}p{4.1cm}}
\toprule
\textbf{Problem} & \textbf{Design response} & \textbf{Required evidence} \\
\midrule
Ambiguous occupation descriptions & Catalogue retrieval and optional
  reranking & Heldout coding accuracy and error analysis \\
Related classification targets & Separate classifiers and explicit
  coherence rules & Coverage, detection errors, and review behaviour \\
Language and dialect variation & Multilingual input processing and
  language-specific answer handling & Per-language evaluation and
  normalisation ablations \\
Conversational survey administration & Controlled interview flow,
  persistence, and HITL review & Functional checks and operational
  measurements; human studies for burden and acceptability \\
\bottomrule
\end{tabular}
\end{table}

\section{Research Approach}
\label{sec:opportunity}

The system combines a finite-state conversation controller with directly
invoked processing modules. CrewAI provides agent abstractions for some
modules, while questionnaire transitions and module invocation remain
under application control. The implementation does not use autonomous
hierarchical delegation. This architecture makes the routing rules and
the points at which classification and review occur inspectable.

Retrieval is studied as a configurable classification component. A flat
configuration searches the unit-group catalogue directly; a hierarchical
configuration restricts successive searches using predicted parent
groups. Catalogue enrichment and embedding size provide additional
experimental factors. Multilingual E5 supplies model sizes with different
resource requirements~\cite{wang2024multilinguale5}. The operational
default and the best-performing benchmark configuration are kept distinct
when interpreting the results.

The evaluation uses three forms of evidence: external labelled occupation
data, synthetic test cases, and software and operational checks. The
external WISCO dataset supports occupation-classification comparisons
\cite{wisco2023}; synthetic cases support narrower studies of classifiers,
coherence rules, and conversation execution. These evidence sources
address different questions and are not pooled as if they represented
the same population or validation setting.

\section{Research Questions}
\label{sec:research-questions}

\begin{enumerate}[label=\textbf{RQ\arabic*.}, leftmargin=3.1em]
  \item \textbf{Retrieval effectiveness.} How do flat and hierarchical
  retrieval, catalogue enrichment, and embedding choice affect ISCO-08
  classification effectiveness on an external heldout benchmark?

  \item \textbf{Classification integration and review.} How can occupation,
  industry, and education coding be integrated with transparent
  cross-standard screening, and what evidence supports the resulting
  classifications and escalation rules?

  \item \textbf{Language robustness.} What do the available language-specific
  evaluations establish about multilingual and dialectal input, and does
  dictionary-based Gulf Arabic normalisation improve classification on
  the tested sample?

  \item \textbf{Operational feasibility.} To what extent can the modular
  conversation system execute its interview paths, validation, persistence,
  and review workflow reliably, and which claims still require evaluation
  with human respondents?
\end{enumerate}

\section{Research Objectives}
\label{sec:objectives}

The aim is to design and evaluate an integrated conversational survey
prototype while establishing the limits of its evidence.
Table~\ref{tab:ch1-objectives} operationalises the research questions.
The measures describe how each objective is assessed, rather than
postulated accuracy or cost targets. Chapter~\ref{chap:results} reports
which assessments were completed and where evidence remains incomplete.

\begin{table}[htbp]
\centering
\caption{Objectives, research questions, and assessment measures.}
\label{tab:ch1-objectives}
\small\renewcommand{\arraystretch}{1.25}
\begin{tabular}{>{\centering\bfseries}p{0.5cm}
                >{\raggedright}p{5.1cm}
                >{\raggedright\arraybackslash}p{6.3cm}}
\toprule
\textbf{ID} & \textbf{Objective} & \textbf{Assessment} \\
\midrule
O1 & Implement a controlled conversational questionnaire (RQ4) &
  Interview-path and transition correctness; response preservation;
  completed and failed synthetic sessions \\
O2 & Compare ISCO-08 retrieval configurations (RQ1) &
  Exact unit-group and higher-level accuracy, uncertainty intervals,
  paired configuration comparisons, and error analysis \\
O3 & Implement ISIC Rev.~4 industry coding (RQ2) &
  Catalogue coverage and class-level performance on the available
  synthetic evaluation; explicit treatment of uncovered classes \\
O4 & Implement education-attainment and field coding (RQ2) &
  Separate evaluation of attainment level and detailed field;
  catalogue coverage and synthetic-data limitations \\
O5 & Implement auditable coherence screening and escalation (RQ2) &
  Rule-level behaviour, available detection-error analysis, and
  verification that HIGH-severity cases enter the review workflow \\
O6 & Assess language and dialect handling (RQ3) &
  Results by the languages represented in each dataset;
  normalisation ablation; disclosure of missing native-speaker validation \\
O7 & Integrate pre-fill and response validation (RQ4) &
  Preservation and correction of stored fields; deterministic
  consistency-rule checks on relevant paths \\
O8 & Support traceable operation and reporting (RQ4) &
  Stored method and review records, report generation, processing time,
  resource constraints, and deployment checks \\
\bottomrule
\end{tabular}
\end{table}

\section{Research Contributions}
\label{sec:contributions}

\noindent\textbf{C1 --- Comparative ISCO-08 retrieval evidence.}
The study compares flat and hierarchical retrieval under documented
configurations on an external benchmark. It also evaluates catalogue
enrichment and embedding choice. Within the evaluated setting, flat
retrieval outperforms the hierarchical configuration. This finding is
task- and configuration-specific; it is not a general rejection of
hierarchical information retrieval.

\noindent\textbf{C2 --- Transparent cross-standard screening.}
The SRE combines hand-built occupation--industry and
occupation--education associations with explicit scoring and severity
rules. It provides an inspectable mechanism for identifying cases for
HITL review. Its scores are heuristic indicators, not empirically
calibrated probabilities of correctness; expert validation and assessment
of legitimate atypical cases remain necessary.

\noindent\textbf{C3 --- An integrated conversational survey prototype.}
The implementation connects questionnaire routing, multilingual input,
multiple classification targets, deterministic validation, persistence,
and review in one workflow. This is a system integration contribution.
It does not establish that these individual techniques are new or that
multi-agent orchestration itself improves accuracy.

\noindent\textbf{C4 --- Evaluation and diagnostic artifacts.}
The project provides configuration-specific benchmark results, synthetic
evaluation cases, functional checks, and operational records. These
artifacts expose coverage limitations and negative results, including the
absence of a measured benefit from Gulf normalisation in the tested
sample. They support inspection and subsequent replication without
substituting synthetic evidence for respondent validation.

\section{Scope and Limitations}
\label{sec:scope}

The scope is a text-based research prototype with UAE-oriented fields,
not a representative national LFS or an approved replacement for an
official collection system. Sampling design, population weighting, and
estimation of national labour-market indicators are outside the study.
The implemented questionnaire paths and rules require further review
before claims of full compliance with official survey standards can be
made.

External occupation evidence comes from WISCO and does not represent
the intended UAE respondent population or natural conversational responses
in the five interface languages.
ISIC and ISCED-F evaluation relies on taxonomy-grounded synthetic data,
and their indexed catalogues have incomplete coverage. Coverage also
varies by catalogue profile because non-standard legacy entries are
excluded from enriched collections. Chapter~\ref{chap:architecture}
reports the relevant inventories, and Chapter~\ref{chap:evaluation}
specifies the denominators and applicability of each evaluation.

The Person Register reuses records from this application's database.
Integration with a national identity or administrative register is
outside the scope. Functional checks of pre-fill do not measure
respondents' willingness to confirm stored answers or the resulting
measurement error. Similarly, classification confidence thresholds and
SRE severity rules have not been established as calibrated measures of
risk on representative respondent data.

A separate synthetic operational run includes 30 scripted English
interviews. It measures system execution and processing time, not human
completion time, satisfaction, or participation. The planned randomised
study with human respondents was not conducted. Accordingly, this thesis
does not establish reductions in respondent burden or collection cost,
non-inferiority to interviewer administration, or population-level
fairness. Ethics approval, recruitment, native-speaker review, independent
coding, and an appropriately justified analysis plan are prerequisites
for those assessments.

\section{Thesis Organisation}
\label{sec:organisation}

Chapter~\ref{chap:relatedwork} reviews occupation coding, retrieval,
multilingual processing, and survey-system integration, then maps the
research gaps to the questions above.
Chapter~\ref{chap:architecture} presents the system architecture,
questionnaire controller, classifiers, and coherence rules.
Chapter~\ref{chap:implementation} describes the implementation and
functional checks.
Chapter~\ref{chap:evaluation} specifies the datasets, comparisons,
statistical methods, and the proposed human-study protocol.
Chapter~\ref{chap:results} reports the measured outcomes and their validity
limits.
Chapter~\ref{chap:conclusion} discusses the contributions, remaining gaps,
and directions for further research.
